%% =========================================================
%% Ambient sources and source-only density reduction.
%% Source: September 24, 2026 polynomial-PEPS preprint, Theorem 5.2.
%% The prescribed memory chain and common readout are explicit hypotheses.
%% Circuit, garbage-storage, gate-participation, and network identifications remain open.
%% =========================================================

\chapter{Ambient Sources and Density Reduction}\label{ch:source_reduction}

Local endpoint maps transport the Gaussian product source into ambient
spaces while preserving its expectation. A separate contraction calculation
bounds the physical density error when selected gates are approximated and
rescaled. The input and output vectors retain their subnormalization.
These calculations give the source and error estimates used in
\cite[Theorem~5.2]{OpenAI2026PolynomialPEPS}.

\section{Ambient endpoint coordinates}

\begin{definition}[Transport through four endpoint maps]
    \label{def:reduction_ambient_transport}
    \lean{QICLean.ComplexGaussian.ambientSchmidtVector}
    \lean{QICLean.ComplexGaussian.ambientSchmidtSource}
    \lean{QICLean.ComplexGaussian.sourceTransport}
    \lean{QICLean.ComplexGaussian.ambientSourceU}
    \lean{QICLean.ComplexGaussian.ambientSourceV}
    \lean{QICLean.ComplexGaussian.ambientSampledSource}
    \lean{QICLean.ComplexGaussian.ambientSourceCorrection}
    \lean{QICLean.ComplexGaussian.ambientSchmidtVector_eq_mulVec}
    \lean{QICLean.ComplexGaussian.sourceTransport_kronecker}
    \lean{QICLean.ComplexGaussian.ambientSourceU_kronecker_ambientSourceV}
    \lean{QICLean.ComplexGaussian.ambientSampledSource_eq_transport}
    \lean{QICLean.ComplexGaussian.sourceTransport_schmidtSource}
    \lean{QICLean.ComplexGaussian.ambientSourceCorrection_eq_transport}
    \leanok
    \uses{def:source_product}
    Let $A,C,W,X,Y,Z$ be finite index sets. Choose arbitrary fixed endpoint
    maps $E:\C^A\to\C^W$, $F:\C^A\to\C^X$,
    $\widetilde E:\C^C\to\C^Y$, and
    $\widetilde F:\C^C\to\C^Z$.
    For nonnegative real weights $\lambda,\mu$, define
    \begin{align}
        |\eta_{E,F}\rangle
                      & =\sum_a\sqrt{\lambda_a}\,E|a\rangle\otimes F|a\rangle,
        \label{eq:reduction_ambient_vector}                                                    \\
        X_{\mathrm{amb}}
                      & =|\eta_{E,F}\rangle\langle\widetilde\eta_{\widetilde E,\widetilde F}|,
        \label{eq:reduction_ambient_target}                                                    \\
        \mathcal T(M) & =(E\otimes F)M(\widetilde E\otimes\widetilde F)^\dagger.
        \label{eq:reduction_ambient_transport}
    \end{align}
    The second ambient vector uses the weights $\mu$.
    For the coordinate sampled factors of Definition~\ref{def:source_product},
    put $U_j^{\mathrm{amb}}=EU_j\widetilde E^\dagger$ and
    $V_j^{\mathrm{amb}}=FV_j\widetilde F^\dagger$. Then
    \begin{align}
        \mathcal T(U\otimes V)
                         & =(EU\widetilde E^\dagger)\otimes(FV\widetilde F^\dagger),
        \label{eq:reduction_ambient_product}                                                    \\
        \widehat X_{\mathrm{amb}}
                         & =\frac1k\sum_{j=0}^{k-1}U_j^{\mathrm{amb}}\otimes V_j^{\mathrm{amb}}
        =\mathcal T(\widehat X),
        \label{eq:reduction_ambient_sample}                                                     \\
        X_{\mathrm{amb}} & =\mathcal T(X),
        \label{eq:reduction_ambient_exact}                                                      \\
        \Delta_{\mathrm{amb}}
                         & =\widehat X_{\mathrm{amb}}-X_{\mathrm{amb}}=\mathcal T(\Delta).
        \label{eq:reduction_ambient_correction}
    \end{align}
    In particular, the represented ambient vector is the endpoint
    transport of its paired coordinate Schmidt vector.
    These are the ambient-coordinate forms of
    \cite[Equation~(5.6)]{OpenAI2026PolynomialPEPS}.
    % Source: eq:compression-random-source, 04-compression.tex:279--309.
    % The algebraic transport identities hold for arbitrary endpoint maps.
\end{definition}

\begin{theorem}[Rank-one expansion of transported sources]
    \label{thm:reduction_transport_rank_one}
    \lean{QICLean.ComplexGaussian.sourceTransport_eq_sum_rankOne}
    \leanok
    \uses{def:reduction_ambient_transport}
    Let $A,C$ be finite index sets and let $W,X,Y,Z$ be arbitrary index sets.
    For endpoint matrices $E,F,\widetilde E,\widetilde F$ of the indicated
    sizes, define the coordinate arrays
    $u_{ab}(w,x)=E_{wa}F_{xb}$ and
    $v_{cd}(y,z)=\widetilde E_{yc}\widetilde F_{zd}$.
    Every matrix $M$ with rows indexed by $A\times A$ and columns indexed
    by $C\times C$ satisfies
    \begin{align}
        \mathcal T(M)
         & =\sum_{a,b\in A}\sum_{c,d\in C} M_{(a,b),(c,d)}\,u_{ab}v_{cd}^{\dagger}.
        \label{eq:reduction_rank_one_expansion}
    \end{align}
    No orthonormality or spanning condition on the endpoint matrices is needed.
\end{theorem}

\begin{proof}\leanok
    \uses{def:reduction_ambient_transport}
    Expand both finite matrix products in
    $\mathcal T(M)=(E\otimes F)M(\widetilde E\otimes\widetilde F)^{\dagger}$.
    The term indexed by $(a,b,c,d)$ at row $(w,x)$ and column $(y,z)$ is
    $M_{(a,b),(c,d)}E_{wa}F_{xb}
        \overline{\widetilde E_{yc}\widetilde F_{zd}}$, which is the corresponding
    entry of the rank-one summand in (\ref{eq:reduction_rank_one_expansion}).
\end{proof}

\begin{theorem}[Unbiased ambient product sources]
    \label{thm:reduction_ambient_unbiased}
    \lean{QICLean.ComplexGaussian.integrable_ambientSourceU}
    \lean{QICLean.ComplexGaussian.integrable_ambientSourceV}
    \lean{QICLean.ComplexGaussian.integrable_ambientSampledSource_entry}
    \lean{QICLean.ComplexGaussian.integral_ambientSampledSource_entry}
    \lean{QICLean.ComplexGaussian.integrable_ambientSampledSource}
    \lean{QICLean.ComplexGaussian.integral_ambientSampledSource}
    \lean{QICLean.ComplexGaussian.integrable_ambientSourceCorrection}
    \lean{QICLean.ComplexGaussian.integral_ambientSourceCorrection}
    \leanok
    \uses{def:reduction_ambient_transport}
    For the actual circular complex Gaussian law, each ambient factor
    $U_j^{\mathrm{amb}},V_j^{\mathrm{amb}}$ is matrix-valued integrable;
    this integrability holds for arbitrary fixed real coefficient weights.
    If $k>0$ and the endpoint weights are nonnegative,
    both $\widehat X_{\mathrm{amb}}$ and $\Delta_{\mathrm{amb}}$ are
    integrable, entrywise and as matrices, and
    \begin{align}
        \mathbb E\widehat X_{\mathrm{amb}} & =X_{\mathrm{amb}},
        \label{eq:reduction_ambient_mean}                       \\
        \mathbb E\Delta_{\mathrm{amb}}     & =0.
        \label{eq:reduction_ambient_centered}
    \end{align}
    These expectation identities hold for arbitrary fixed endpoint maps;
    the isometry equations govern the separate norm-preservation estimate
    of Theorem~\ref{thm:source_physical_reinsertion}.
    This is the unbiased replacement in
    \cite[Equation~(5.6)]{OpenAI2026PolynomialPEPS}.
    % Source: eq:compression-random-source, 04-compression.tex:294--309.
    % Abstract Schmidt existence is available; the explicit endpoint maps are data here.
\end{theorem}

\begin{proof}\leanok
    \uses{thm:source_unbiased}
    Each factor entry is a finite linear combination of integrable
    Gaussian coordinates. Finite matrix evaluation gives factor
    integrability. The transport identities express every ambient source
    entry as a finite linear combination of coordinate source entries.
    Integrate these sums and use $\mathbb E\widehat X=X$ from
    Theorem~\ref{thm:source_unbiased}. Transport of the target gives
    (\ref{eq:reduction_ambient_mean}); subtracting it gives
    (\ref{eq:reduction_ambient_centered}). Finite evaluation combines the
    entry identities into the matrix-valued identities.
\end{proof}

\section{Subnormalized vectors and changing memories}

\begin{theorem}[Absolute density error for subnormalized vectors]
    \label{thm:reduction_pure_error}
    \lean{Matrix.euclideanOuterProduct}
    \lean{Matrix.trace_euclideanOuterProduct_conjTranspose_mul}
    \lean{Matrix.rectangularTraceNorm_euclideanOuterProduct_le}
    \lean{Matrix.euclideanOuterProduct_self_sub}
    \lean{Matrix.rectangularTraceNorm_pure_density_sub_le}
    \lean{Matrix.rectangularTraceNorm_pure_density_sub_le_two}
    \lean{Matrix.rectangularTraceNorm_partialTraceRight_pure_density_sub_le}
    \lean{Matrix.rectangularTraceNorm_partialTraceRight_pure_density_sub_le_two}
    \leanok
    For finite Euclidean spaces, the rectangular outer product has entries
    $(|u\rangle\langle v|)_{ij}=u_i\overline{v_j}$ and satisfies
    $\lVert|u\rangle\langle v|\rVert_1\leq\lVert u\rVert\lVert v\rVert$.
    For vectors $v,w$ in the same space,
    \begin{align}
        \lVert|v\rangle\langle v|-|w\rangle\langle w|\rVert_1
         & \leq(\lVert v\rVert+\lVert w\rVert)\lVert v-w\rVert.
        \label{eq:reduction_pure_error}
    \end{align}
    Discarding a finite register preserves this upper bound. In particular,
    if $\lVert v\rVert,\lVert w\rVert\leq1$, both the full and retained
    density errors are at most $2\lVert v-w\rVert$.
    Zero vectors and zero-dimensional spaces are allowed.
    This is the vector-to-density estimate preceding
    \cite[Equation~(5.3)]{OpenAI2026PolynomialPEPS}.
    % Source: eq:compression-effect-circuit-error, 04-compression.tex:218--229.
\end{theorem}

\begin{proof}\leanok
    \uses{thm:random_trace_duality, thm:source_branch_sum,
        thm:random_exterior_contraction}
    The trace pairing is
    $\tr((|u\rangle\langle v|)^\dagger U)=\langle u,Uv\rangle$.
    A test contraction $U$ and Cauchy--Schwarz bound its absolute value by
    $\lVert u\rVert\lVert v\rVert$, so trace-norm duality gives the
    outer-product bound. Split the density difference as
    $|v-w\rangle\langle v|+|w\rangle\langle v-w|$ and use the triangle
    inequality to obtain (\ref{eq:reduction_pure_error}). Partial-trace
    contraction gives the retained-density bound. The two upper norm
    bounds give the factor two.
\end{proof}

\begin{theorem}[Telescoping and rescaling through changing memory spaces]
    \label{thm:reduction_chain}
    \lean{Matrix.contractionPrefix}
    \lean{Matrix.l2_opNorm_one_le}
    \lean{Matrix.contractionPrefix_norm_le_one}
    \lean{Matrix.contractionPrefix_sub_norm_le}
    \lean{Matrix.contractionPrefix_sub_norm_le_uniform}
    \lean{Matrix.contractionPrefix_sub_norm_le_supported}
    \lean{NormedSpace.rescale_approximation}
    \leanok
    Let $D_t$ be finite index sets, and let
    $G_t,H_t:\C^{D_t}\to\C^{D_{t+1}}$ be contractions for every
    $t\in\N$. Define the chronological prefixes by
    $\Pi_G(0)=\Id$ and $\Pi_G(t+1)=G_t\Pi_G(t)$.
    Each prefix is a contraction, including the identity on an empty
    space, and
    \begin{align}
        \lVert\Pi_G(n)-\Pi_H(n)\rVert
         & \leq\sum_{t=0}^{n-1}\lVert G_t-H_t\rVert.
        \label{eq:reduction_telescoping}
    \end{align}
    A uniform error bound $\lVert G_t-H_t\rVert\leq\delta$ gives the
    bound $n\delta$. More generally, for a finite marked set $S\subset\N$
    and $\delta\geq0$, errors at most $\delta$ on $S$ and exact equality
    outside $S$ give $\lVert\Pi_G(n)-\Pi_H(n)\rVert\leq|S|\delta$.
    In a real normed space, if $\lVert G\rVert\leq1$ and
    $\lVert A-G\rVert\leq\delta$ with $\delta\geq0$, then
    $\lVert A/(1+\delta)\rVert\leq1$ and
    $\lVert A/(1+\delta)-G\rVert\leq2\delta$.
    These are the contraction estimates used before
    \cite[Equation~(5.3)]{OpenAI2026PolynomialPEPS}.
    % Source: eq:compression-effect-circuit-error, 04-compression.tex:210--229.
    % Gate families are indexed by all naturals; finite circuits can be extended afterward.
\end{theorem}

\begin{proof}\leanok
    Induct on the prefix length. At the next stage, use
    \begin{align}
        \Pi_G(t+1)-\Pi_H(t+1)
         & =G_t(\Pi_G(t)-\Pi_H(t))+(G_t-H_t)\Pi_H(t).
        \label{eq:reduction_telescoping_step}
    \end{align}
    Submultiplicativity and the prefix contraction bound add the next
    gate error. If the gates agree outside $S$, only indices in
    $S\cap\{0,\ldots,n-1\}$ contribute; their number is at most $|S|$.
    For rescaling, $\lVert A\rVert\leq1+\delta$ gives the contraction
    bound. The decomposition
    $A/(1+\delta)-G=(A-G)/(1+\delta)
        +((1+\delta)^{-1}-1)G$ gives the error bound $2\delta$.
\end{proof}

\section{Actual readout densities and marked error budgets}

\begin{theorem}[Physical densities from a common contractive readout]
    \label{thm:reduction_readout}
    \lean{Matrix.sourceOnlyReadoutVector}
    \lean{Matrix.sourceOnlyReadoutDensity}
    \lean{Matrix.sourceOnlyReadoutVector_norm_le_one}
    \lean{Matrix.sourceOnlyReadoutVector_sub_norm_le}
    \lean{Matrix.rectangularTraceNorm_sourceOnlyReadoutDensity_sub_le_prefix}
    \lean{Matrix.rectangularTraceNorm_sourceOnlyReadoutDensity_sub_le_sum}
    \lean{Matrix.sourceOnlyReadoutDensity_zero}
    \lean{Matrix.rectangularTraceNorm_sourceOnlyReadoutDensity_sub_zero}
    \leanok
    Let $\psi\in\C^{D_0}$ satisfy $\lVert\psi\rVert\leq1$, and let
    $K:\C^{D_n}\to\C^{P\times E}$ be a common contraction into finite
    physical and discarded registers. Define the actual final vector and
    retained density by
    \begin{align}
        v_G    & =K\Pi_G(n)\psi,
        \label{eq:reduction_readout_vector}       \\
        \rho_G & =\Tr_E(|v_G\rangle\langle v_G|).
        \label{eq:reduction_readout_density}
    \end{align}
    For arbitrary gate families,
    $\lVert v_H-v_G\rVert\leq\lVert\Pi_H(n)-\Pi_G(n)\rVert$.
    If both families are contractions, $v_G,v_H$ have norm at most one and
    \begin{align}
        \lVert\rho_H-\rho_G\rVert_1
         & \leq2\lVert\Pi_H(n)-\Pi_G(n)\rVert
        \leq2\sum_{t=0}^{n-1}\lVert H_t-G_t\rVert.
        \label{eq:reduction_readout_error}
    \end{align}
    For $n=0$ the densities coincide and their trace-norm difference is
    zero; this equality holds for arbitrary $\psi,K,G,H$.
    These are the actual vector and density estimates used in
    \cite[Equation~(5.3)]{OpenAI2026PolynomialPEPS}.
    % Source: eq:compression-effect-circuit-error, 04-compression.tex:214--229.
    % A caller supplies the common readout and its circuit interpretation.
\end{theorem}

\begin{proof}\leanok
    \uses{thm:reduction_chain, thm:reduction_pure_error}
    The common readout gives
    $v_H-v_G=K(\Pi_H(n)-\Pi_G(n))\psi$; its contraction bound and
    $\lVert\psi\rVert\leq1$ give the vector estimate. Prefix and readout
    contraction also give the two upper vector norm bounds. Apply
    Theorem~\ref{thm:reduction_pure_error}, then
    (\ref{eq:reduction_telescoping}). At zero length both prefixes are
    the same identity, so the actual densities agree.
\end{proof}

\begin{theorem}[Half-accuracy density budget for marked approximate gates]
    \label{thm:reduction_marked_budget}
    \lean{Matrix.rescaledGateChain}
    \lean{Matrix.markedRescaledGateChain}
    \lean{Matrix.rectangularTraceNorm_sourceOnlyReadoutDensity_sub_le}
    \lean{Matrix.rectangularTraceNorm_sourceOnlyReadoutDensity_sub_le_half}
    \lean{Matrix.rectangularTraceNorm_sourceOnlyReadoutDensity_marked_sub_le}
    \lean{Matrix.rectangularTraceNorm_sourceOnlyReadoutDensity_marked_sub_le_half}
    \lean{Matrix.rectangularTraceNorm_sourceOnlyReadoutDensity_marked_empty}
    \leanok
    Let $G_t$ be contractions on the prescribed changing memory spaces,
    and let $\psi,K$ satisfy the hypotheses of
    Theorem~\ref{thm:reduction_readout}. Choose a finite marked set $S$
    and $\delta\geq0$, with $\lVert A_t-G_t\rVert\leq\delta$ for
    every $t\in S$. Set
    \begin{align}
        H_t & =\begin{cases}
                   A_t/(1+\delta), & t\in S,    \\
                   G_t,            & t\notin S.
               \end{cases}
        \label{eq:reduction_marked_gates}
    \end{align}
    For the actual densities of (\ref{eq:reduction_readout_density}),
    \begin{align}
        \lVert\rho_H-\rho_G\rVert_1 & \leq4|S|\delta.
        \label{eq:reduction_marked_error}
    \end{align}
    If $|S|>0$ and $\epsilon\geq0$, the choice
    $\delta=\epsilon/(8|S|)$ gives error at most $\epsilon/2$.
    For $S=\varnothing$, $H=G$ and the density error is exactly zero,
    for arbitrary input, readout, and rescaling parameter.
    Unmarked private stages remain exact.
    If all stages are approximated with a uniform error bound and
    $H_t=A_t/(1+\delta)$ for every $t$, the corresponding estimate is
    $4n\delta$; for $n>0$ and $\epsilon\geq0$, taking
    $\delta=\epsilon/(8n)$ gives the same
    half-accuracy bound.
    With the nonprivate gate count $M=|S|$, this is
    \cite[Equation~(5.3)]{OpenAI2026PolynomialPEPS}.
    % Source: eq:compression-effect-circuit-error, 04-compression.tex:199--229.
    % Identifying S, the memory chain, and common K with the concrete circuit is separate.
    % These entries do not construct appended-garbage semantics or a final operator network.
\end{theorem}

\begin{proof}\leanok
    \uses{thm:reduction_chain, thm:reduction_readout}
    Rescaling gives $\lVert H_t\rVert\leq1$ and
    $\lVert H_t-G_t\rVert\leq2\delta$ on $S$; outside $S$ the error
    is zero. The supported telescoping bound gives prefix error at most
    $2|S|\delta$. Formula (\ref{eq:reduction_readout_error}) doubles this
    to (\ref{eq:reduction_marked_error}). For positive $|S|$,
    $4|S|\epsilon/(8|S|)=\epsilon/2$.
    The uniform case uses the uniform telescoping estimate in the same
    way. An empty marked set makes the two chains identical.
\end{proof}
